%========================================
\chapter{数据链路层}
%========================================
\section{基本概念}
\begin{itemize}[leftmargin=2em,itemsep=0pt]
    \item \textbf{链路}：一个节点到相邻节点的物理线路
    \item \textbf{数据链路}：链路 + 通信协议（硬件+软件）
    \item \textbf{帧}：数据链路层的传输单位
\end{itemize}

\subsection{三个核心问题}
\begin{center}
\begin{tikzpicture}[node distance=0cm,>=stealth]
\node[draw,minimum width=1.2cm,minimum height=0.7cm,fill=gray!30] (head) {\small 帧头};
\node[draw,minimum width=3.5cm,minimum height=0.7cm,right=of head,fill=white] (data) {\small 数据（来自网络层）};
\node[draw,minimum width=1.2cm,minimum height=0.7cm,right=of data,fill=gray!30] (tail) {\small 帧尾};
\draw[decorate,decoration={brace,amplitude=4pt,raise=2pt}] (head.north west) -- (tail.north east) node[midway,above=6pt] {\small 封装成帧};
\node[below=0.2cm of head] {\footnotesize 帧定界};
\node[below=0.2cm of tail] {\footnotesize FCS校验};
\end{tikzpicture}
\end{center}

\begin{enumerate}[leftmargin=2em,itemsep=0pt]
    \item \textbf{封装成帧}：添加帧头帧尾，实现帧定界
    \item \textbf{透明传输}：数据中出现帧定界符时的处理
    \begin{itemize}[itemsep=0pt]
        \item \textbf{字节填充}：在数据中的特殊字符前插入转义字符（如PPP协议）
        \item \textbf{比特填充}：又称“零比特填充”，发送端在连续5个“1”后填入一个“0”，接收端在连续5个“1”后删掉紧接着的一个“0”。这确保了数据中不会出现帧定界符（通常是连续6个1）。
    \end{itemize}
    \item \textbf{差错检测}：通过FCS检测比特错误（不纠错）
\end{enumerate}

\section{差错检测}
\subsection{奇偶校验}
在数据后添加1位校验位，使整个数据中"1"的个数为奇数或偶数。

\textbf{局限性}：只能检测奇数位错误，无法检测偶数位错误。

\subsection{CRC循环冗余校验}
CRC是一种强大的差错检测方法，基于多项式除法原理。

\textbf{CRC计算步骤}：
\begin{enumerate}[leftmargin=2em,itemsep=0pt]
    \item 确定生成多项式$G(x)$，将其转换为二进制除数
    \item 在数据后补$n$个0（$n$为生成多项式最高次数）
    \item 用模2除法（异或运算）计算余数
    \item 将余数附加到原数据后发送
\end{enumerate}

\begin{examplebox}[例题3.1 CRC计算]
\textbf{题目}：生成多项式$G(x)=x^3+x^2+1$，数据为\texttt{101001}，求发送的比特序列。
\tcblower
\textbf{步骤}：
\begin{enumerate}[itemsep=0pt]
    \item 除数：$G(x)$对应\texttt{1101}
    \item 被除数：数据后补3个0 $\rightarrow$ \texttt{101001000}
    \item 模2除法得余数：\texttt{001}
\end{enumerate}
\textbf{发送序列}：\texttt{101001} + \texttt{001} = \texttt{101001001}
\end{examplebox}

\section{可靠传输协议}
\subsection{停止-等待协议(SW)}
发送方每发送一个帧后必须等待接收方的确认。

\textbf{关键机制}：
\begin{itemize}[leftmargin=2em,itemsep=0pt]
    \item \textbf{确认(ACK)}：接收方收到正确帧后发送确认
    \item \textbf{超时重传}：发送方超时未收到确认则重传
    \item \textbf{帧编号}：用0/1交替编号，防止重复接收
\end{itemize}

\begin{formulabox}
\textbf{信道利用率}：$U = \frac{T_D}{T_D + RTT + T_A} \approx \frac{T_D}{T_D + RTT}$
\end{formulabox}

\begin{examplebox}[例题3.2 停等协议效率]
\textbf{题目}：帧长1000字节，速率1Mbps，单向传播时延250ms，求信道利用率。
\tcblower
发送时延：$T_D = \frac{1000 \times 8}{10^6} = 8$ms

往返时间：$RTT = 2 \times 250 = 500$ms

信道利用率：$U = \frac{8}{8+500} = \frac{8}{508} \approx$ \textbf{1.57\%}
\end{examplebox}

\subsection{回退N帧协议(GBN)}
允许发送方连续发送多个帧，提高信道利用率。

\textbf{核心特点}：
\begin{itemize}[leftmargin=2em,itemsep=0pt]
    \item 发送窗口$W_T$：$1 < W_T \leq 2^n - 1$（n为序号位数）
    \item 接收窗口$W_R = 1$（只能按序接收）
    \item \textbf{累积确认}：ACKn表示n及之前的帧都已收到
    \item \textbf{出错处理}：重传出错帧及其后所有帧
\end{itemize}

\subsection{选择重传协议(SR)}
只重传出错的帧，效率更高。

\textbf{核心特点}：
\begin{itemize}[leftmargin=2em,itemsep=0pt]
    \item 发送窗口$W_T$：$1 < W_T \leq 2^{n-1}$
    \item 接收窗口$W_R$：$1 < W_R \leq W_T$
    \item \textbf{逐个确认}：只确认正确收到的帧
    \item \textbf{出错处理}：只重传出错的帧
\end{itemize}

\begin{examplebox}[例题3.3 窗口大小计算]
\textbf{题目}：帧序号用3比特表示，GBN和SR协议的发送窗口最大分别为多少？
\tcblower
序号范围：$0 \sim 2^3 - 1 = 0 \sim 7$，共8个序号

\textbf{GBN协议}：$W_T \leq 2^n - 1 = 2^3 - 1 = 7$

\textbf{SR协议}：$W_T \leq 2^{n-1} = 2^2 = 4$
\end{examplebox}

\section{PPP协议}
PPP是目前使用最广泛的点对点数据链路层协议。

\subsection{帧格式}
\begin{center}
\begin{tikzpicture}[node distance=0cm]
\node[draw,minimum width=1cm,minimum height=0.8cm,fill=gray!30] (f1) {F};
\node[draw,minimum width=1cm,minimum height=0.8cm,right=of f1] (a) {A};
\node[draw,minimum width=1cm,minimum height=0.8cm,right=of a] (c) {C};
\node[draw,minimum width=1.5cm,minimum height=0.8cm,right=of c] (p) {协议};
\node[draw,minimum width=3cm,minimum height=0.8cm,right=of p,fill=yellow!20] (data) {数据};
\node[draw,minimum width=1.5cm,minimum height=0.8cm,right=of data] (fcs) {FCS};
\node[draw,minimum width=1cm,minimum height=0.8cm,right=of fcs,fill=gray!30] (f2) {F};
\node[below=0.1cm of f1] {\tiny 7E};
\node[below=0.1cm of a] {\tiny FF};
\node[below=0.1cm of c] {\tiny 03};
\end{tikzpicture}
\end{center}

\subsection{透明传输}
\textbf{异步传输（字节填充）}：
\begin{itemize}[leftmargin=2em,itemsep=0pt]
    \item \texttt{7E}$\rightarrow$\texttt{7D 5E}
    \item \texttt{7D}$\rightarrow$\texttt{7D 5D}
\end{itemize}

\section{ALOHA协议}
ALOHA是最早的随机访问协议。
\begin{itemize}[leftmargin=2em,itemsep=0pt]
    \item \textbf{纯ALOHA}：不监听信道，想发就发。如果冲突（未收到确认），等待随机时间重发。利用率极低（最高约18.4\%）。
    \item \textbf{分隙ALOHA (Slotted ALOHA)}：将时间划分为等长的时隙，只能在时隙开始处发送。最高利用率提升至约36.8\%。
\end{itemize}

\section{PPPoE协议}
PPPoE (PPP over Ethernet) 是在以太网上运行 PPP 协议的技术。
\begin{itemize}[leftmargin=2em,itemsep=0pt]
    \item \textbf{作用}：通过以太网将多个主机连接到宽带接入服务。它结合了以太网的广播能力和 PPP 的身份验证、计费功能。
    \item \textbf{应用}：家庭宽带拨号上网（ADSL/光纤拨号）。
\end{itemize}

\section{CSMA/CD协议}
以太网使用的介质访问控制协议。

\subsection{工作原理}
\begin{center}
\fbox{\textbf{先听后说、边说边听、冲突停止、随机重发}}
\end{center}

\begin{formulabox}
\textbf{争用期}（碰撞窗口）：$2\tau$（往返传播时延）

\textbf{最小帧长} = 争用期 $\times$ 数据率 = $\frac{2d}{v} \times R$

10Mbps以太网：争用期=51.2$\mu$s，最小帧长=64字节
\end{formulabox}

\begin{examplebox}[例题3.4 CSMA/CD计算]
\textbf{题目}：以太网速率10Mbps，最大距离2500m，传播速度$2\times10^8$m/s，求争用期和最小帧长。
\tcblower
单向传播时延：$\tau = \frac{2500}{2\times10^8} = 12.5\mu$s

争用期：$2\tau = 25\mu$s

最小帧长：$L = 2\tau \times R = 25\times10^{-6} \times 10\times10^6 = 250$bit = \textbf{32字节}
\end{examplebox}

\subsection{截断二进制指数退避算法}
碰撞后等待时间 = 争用期 $\times$ 随机数$r$，其中$r \in \{0,1,...,2^k-1\}$，$k=\min(\text{重传次数},10)$

\begin{examplebox}[综合题3.4b ARP协议分析]
\textbf{题目}：观察捕获到的 ARP 请求报文，分析各字段含义，并简述 ARP 协议的工作原理。
\begin{lstlisting}[basicstyle=\ttfamily\footnotesize]
Address Resolution Protocol (request)
    Hardware type: Ethernet (1)
    Protocol type: IPv4 (0x0800)
    Hardware size: 6
    Protocol size: 4
    Opcode: request (1)
    Sender MAC address: 00:0c:29:ab:cd:ef
    Sender IP address: 192.168.1.1
    Target MAC address: 00:00:00:00:00:00 (unknown)
    Target IP address: 192.168.1.2
\end{lstlisting}
\tcblower
\textbf{字段分析}：
\begin{itemize}[itemsep=0pt]
    \item \textbf{Hardware/Protocol type}：标识底层链路（以太网）和网络层协议（IPv4）。
    \item \textbf{Size}：MAC地址长6字节，IP地址长4字节。
    \item \textbf{Opcode}：1 表示请求 (Request)，2 表示响应 (Reply)。
    \item \textbf{Sender/Target}：包含发送端的 IP/MAC 以及目的端的 IP。请求时目的 MAC 为全0。
\end{itemize}

\textbf{ARP工作原理}：
\begin{enumerate}[itemsep=0pt]
    \item \textbf{查询缓存}：主机先在自己的 ARP 缓存表中查找目的 IP 对应的 MAC。
    \item \textbf{广播请求}：若找不到，则发送一个 \textbf{ARP请求帧}（目的MAC为广播地址 \texttt{FF-FF...}）。
    \item \textbf{单播响应}：目的主机匹配到自己的 IP 后，回复一个 \textbf{ARP响应帧}（单播给发送方），告诉其自己的 MAC 并更新自己的缓存。
    \item \textbf{更新缓存}：源主机收到后，将目的 IP 和 MAC 存入缓存。
\end{enumerate}
\end{examplebox}

\section{以太网标准 (IEEE 802.3)}
以太网标准采用特定的命名规范来描述其传输速率、信号方式和传输介质。

\textbf{命名格式}：\fbox{\texttt{<速率> Base - <介质>}}

\begin{itemize}[leftmargin=2em,itemsep=0pt]
    \item \textbf{速率}：数据率，单位为 Mbps。如 10、100、1000。
    \item \textbf{Base}：指基带传输（Baseband），即信号直接在物理介质上传输。
    \item \textbf{介质}：
    \begin{itemize}[itemsep=0pt]
        \item \texttt{5}：粗同轴电缆，最大段长 500m。
        \item \texttt{2}：细同轴电缆，最大段长 185m (约 200m)。
        \item \texttt{T}：双绞线（Twisted Pair），目前最常用。
        \item \texttt{F}：光纤（Fiber）。
    \end{itemize}
\end{itemize}

\textbf{常见标准}：
\begin{table}[htbp]
\centering\small
\begin{tabular}{llll}
\toprule
\textbf{标准名称} & \textbf{速率} & \textbf{介质} & \textbf{最大段长} \\
\midrule
10Base-5 & 10 Mbps & 粗同轴电缆 & 500m \\
10Base-2 & 10 Mbps & 细同轴电缆 & 185m \\
10Base-T & 10 Mbps & 双绞线(UTP) & 100m \\
100Base-T & 100 Mbps & 双绞线 & 100m (快速以太网) \\
1000Base-SX & 1000 Mbps & 短波多模光纤 & 550m (吉比特以太网) \\
\bottomrule
\end{tabular}
\end{table}

\section{MAC地址与ARP}
\subsection{MAC地址格式}
\begin{itemize}[leftmargin=2em,itemsep=0pt]
    \item 48位（6字节），固化在网卡的ROM中
    \item 格式：前3字节OUI（厂商标识）+ 后3字节设备标识
    \item 广播地址：\texttt{FF-FF-FF-FF-FF-FF}
\end{itemize}

\subsection{ARP协议}
ARP用于将IP地址解析为MAC地址。

\textbf{工作过程}：
\begin{enumerate}[leftmargin=2em,itemsep=0pt]
    \item 检查本地ARP缓存
    \item 若没有，广播ARP请求
    \item 目标主机单播回复MAC地址
    \item 将映射存入ARP缓存
\end{enumerate}

\section{以太网交换机}
\subsection{集线器vs交换机}
\begin{table}[htbp]
\centering\small
\begin{tabular}{lcc}
\toprule
\textbf{特性} & \textbf{集线器} & \textbf{交换机} \\
\midrule
工作层级 & 物理层 & 数据链路层 \\
转发依据 & 无（广播） & MAC地址 \\
冲突域 & 所有端口共享 & 每端口独立 \\
通信方式 & 半双工 & 全双工 \\
\bottomrule
\end{tabular}
\end{table}

\subsection{自学习与转发}
\textbf{转发规则}：
\begin{enumerate}[leftmargin=2em,itemsep=0pt]
    \item \textbf{自学习}：记录源MAC与入端口的映射
    \item \textbf{明确转发}：表中有目的MAC，转发到对应端口
    \item \textbf{泛洪}：表中无目的MAC，转发到除入端口外所有端口
    \item \textbf{过滤}：目的MAC对应入端口，直接丢弃
\end{enumerate}

\begin{examplebox}[例题3.5 交换机转发]
\textbf{题目}：交换机初始表为空，依次发生：(1)A$\to$B (2)B$\to$A (3)C$\to$D，描述操作。
\tcblower
\begin{tabular}{|c|l|l|l|}
\hline
\textbf{通信} & \textbf{自学习} & \textbf{转发} & \textbf{帧交换表} \\
\hline
A$\to$B & 登记A-端口1 & 泛洪(2,3,4) & A:1 \\
B$\to$A & 登记B-端口2 & 明确转发(1) & A:1, B:2 \\
C$\to$D & 登记C-端口3 & 泛洪(1,2,4) & A:1, B:2, C:3 \\
\hline
\end{tabular}
\end{examplebox}

\section{生成树协议(STP)}
\subsection{问题背景}
冗余链路提高可靠性，但会形成环路，导致：
\begin{itemize}[leftmargin=2em,itemsep=0pt]
    \item \textbf{广播风暴}：广播帧在环路中无限循环
    \item \textbf{帧重复}：主机收到多份相同帧
    \item \textbf{MAC表震荡}：地址映射不断变化
\end{itemize}

\subsection{STP关键概念}
\begin{enumerate}[leftmargin=2em,itemsep=0pt]
    \item \textbf{根桥(Root Bridge)}：网络中优先级最高的交换机
    \item \textbf{根端口(RP)}：非根桥上到根桥路径最短的端口
    \item \textbf{指定端口(DP)}：每个网段上到根桥路径最短的端口
    \item \textbf{阻塞端口}：既非根端口也非指定端口，被阻塞
\end{enumerate}

\section{虚拟局域网(VLAN)}
VLAN将物理局域网在逻辑上划分为多个广播域。

\subsection{VLAN的作用}
\begin{itemize}[leftmargin=2em,itemsep=0pt]
    \item \textbf{隔离广播域}：减少广播风暴的影响范围
    \item \textbf{提高安全性}：不同VLAN之间默认隔离
    \item \textbf{灵活管理}：按业务或部门划分
\end{itemize}

\subsection{802.1Q帧格式}
在标准以太网帧中插入4字节VLAN标记：
\begin{itemize}[leftmargin=2em,itemsep=0pt]
    \item \textbf{TPID}（2字节）：固定为\texttt{0x8100}
    \item \textbf{TCI}（2字节）：包含VID（12位VLAN标识符，范围1-4094）
\end{itemize}

\subsection{端口类型}
\begin{itemize}[leftmargin=2em,itemsep=0pt]
    \item \textbf{Access端口}：连接终端设备，只属于一个VLAN
    \begin{itemize}[itemsep=0pt]
        \item 接收：仅收未打标签的帧，按端口PVID打标签
        \item 发送：仅转发VID与端口PVID一致的帧，并去标签
    \end{itemize}
    \item \textbf{Trunk端口}：连接交换机，可传输多个VLAN的帧
    \begin{itemize}[itemsep=0pt]
        \item 接收：未打标签的帧按PVID打标签，已打标签的帧保留标签
        \item 发送：VID与PVID一致的帧去标签转发，否则带标签转发
    \end{itemize}
\end{itemize}

\begin{examplebox}[例题3.6 VLAN通信]
\textbf{题目}：交换机S1的端口1-4分别连接主机A、B、C、D，端口1、2属于VLAN10，端口3、4属于VLAN20。A发送广播帧，哪些主机能收到？
\tcblower
广播帧只在同一VLAN内传播。

A属于VLAN10，因此广播帧只会发送到同属VLAN10的端口。

\textbf{答案}：只有B能收到（B也在VLAN10），C和D在VLAN20，收不到。
\end{examplebox}

\section{数据链路层综合例题}

\begin{examplebox}[例题3.7 多交换机自学习（综合题）]
\textbf{题目}：如图所示网络拓扑，主机A-H的MAC地址与主机名相同。交换机1的端口1-4分别连接A、B、交换机2、C；交换机2的端口1-4分别连接交换机1、D、E、交换机3；交换机3的端口1-4分别连接F、交换机2、G、H。依次发生通信：(1)B→C (2)D→A (3)G→D (4)E→H (5)C→B (6)F→G，给出各交换机的自学习过程和最终帧交换表。
\tcblower
\textbf{通信过程分析}：

\begin{center}\small
\begin{tabular}{|c|l|l|l|}
\hline
\textbf{通信} & \textbf{交换机1} & \textbf{交换机2} & \textbf{交换机3} \\
\hline
B→C & 登记B:2，泛洪 & 登记B:1，泛洪 & 登记B:2，泛洪 \\
D→A & 登记D:3，泛洪 & 登记D:2，泛洪 & 登记D:2，泛洪 \\
G→D & 收不到 & 登记G:4，明确转发 & 登记G:3，明确转发 \\
E→H & 登记E:3，泛洪 & 登记E:3，泛洪 & 登记E:2，泛洪 \\
C→B & 登记C:4，明确转发 & 收不到 & 收不到 \\
F→G & 收不到 & 收不到 & 登记F:1，明确转发 \\
\hline
\end{tabular}
\end{center}

\textbf{最终帧交换表}：
\begin{center}\small
\begin{tabular}{|c|c||c|c||c|c|}
\hline
\multicolumn{2}{|c||}{\textbf{交换机1}} & \multicolumn{2}{c||}{\textbf{交换机2}} & \multicolumn{2}{c|}{\textbf{交换机3}} \\
\hline
MAC & 端口 & MAC & 端口 & MAC & 端口 \\
\hline
B & 2 & B & 1 & B & 2 \\
D & 3 & D & 2 & D & 2 \\
E & 3 & G & 4 & G & 3 \\
C & 4 & E & 3 & E & 2 \\
  &   &   &   & F & 1 \\
\hline
\end{tabular}
\end{center}
\end{examplebox}

\begin{examplebox}[例题3.8 2014年真题]
\textbf{题目}：某以太网拓扑及交换机当前转发表如下图所示，主机00-e1-d5-00-23-a1向主机00-e1-d5-00-23-c1发送1个数据帧，主机00-e1-d5-00-23-c1收到该帧后，向主机00-e1-d5-00-23-a1发送1个确认帧，交换机对这两个帧的转发端口分别是？

\begin{center}\small
\begin{tabular}{|c|c|}
\hline
目的地址 & 端口 \\
\hline
00-e1-d5-00-23-b1 & 2 \\
00-e1-d5-00-23-a1 & 1 \\
\hline
\end{tabular}
\end{center}
（交换机有3个端口，端口1连a1，端口2连b1，端口3连c1）

A. \{3\}和\{1\} \quad B. \{2,3\}和\{1\} \quad C. \{2,3\}和\{1,2\} \quad D. \{1,2,3\}和\{1\}
\tcblower
\textbf{分析}：

(1) a1→c1：查表，c1不在表中，执行\textbf{泛洪}，转发到端口\{2,3\}（除入端口1外所有端口）

(2) c1→a1：查表，a1在端口1，执行\textbf{明确转发}，转发到端口\{1\}

\textbf{答案}：B
\end{examplebox}

\section{本章习题}

\begin{examplebox}[填空题]
\textbf{1.} 流量控制是指\underline{使得发送端的发送速率不要太快，接收端来得及接收}。

\textbf{2.} 允许发送站发送多个帧而不需要马上应答，这就是\underline{滑动窗口}协议。

\textbf{3.} HDLC有三种不同类型的帧，分别称为\underline{信息帧}、\underline{监督帧}、\underline{无编号帧}。

\textbf{4.} PPP协议包括三个部分：\underline{PPP}、\underline{NCP}和\underline{LCP}。

\textbf{5.} 使用RARP将\underline{物理地址}转换为\underline{IP地址}。（注：ARP是将IP地址转换为物理地址）
\end{examplebox}

\begin{examplebox}[选择题]
\textbf{1.} 局域网中访问冲突的根源是（\textbf{B}）\\
A. 独占介质 \quad B. 共享介质 \quad C. 引入MAC子层 \quad D. 规则的拓扑结构

\textbf{2.} 使用双绞线的10M以太网为（\textbf{C}）\\
A. 10Base-5 \quad B. 10Base-2 \quad C. 10Base-T \quad D. 10Base-F

\textbf{3.} 使用双绞线的100M以太网为（\textbf{C}）\\
A. 100Base-5 \quad B. 100Base-2 \quad C. 100Base-T \quad D. 100Base-F

\textbf{4.} 属于数据链路层的联网设备有（\textbf{C}）\\
A. 中继器 \quad B. 路由器 \quad C. 交换机 \quad D. 网关

\textbf{5.} 在Go-Back-N协议中，当序号为3比特，发送窗口的最大尺寸为（\textbf{C}）\\
A. 5 \quad B. 6 \quad C. 7 \quad D. 8

\textbf{6.} 采用比特同步的帧数据段中出现位串011111001，则比特填充后输出是（\textbf{A}）\\
A. 0111110001 \quad B. 0011111001 \quad C. 0111111001 \quad D. 01111100010

\textbf{7.} 若数据链路的发送窗口尺寸W=5，在发送3帧、并接到2帧的确认帧后，发送方还可连续发送（\textbf{C}）帧\\
A. 2 \quad B. 3 \quad C. 4 \quad D. 1
\end{examplebox}

\begin{examplebox}[判断题]
\textbf{1.} （×）使用RARP将物理地址转换为IP地址。\\
\textit{解析}：RARP是将物理地址转换为IP地址，ARP是将IP地址转换为物理地址。

\textbf{2.} （×）单模光纤比多模光纤细，所以传输速率低。\\
\textit{解析}：单模光纤传输速率更高，传输距离更远。

\textbf{3.} （×）在OSI/RM模型中，网络层传送的是帧。\\
\textit{解析}：网络层传送的是分组/数据报，数据链路层传送的是帧。

\textbf{4.} （√）PPP协议包括三个部分：PPP、NCP和LCP。

\textbf{5.} （√）HDLC是面向比特的通信协议。
\end{examplebox}

\begin{examplebox}[计算题3.9 CSMA/CD]
\textbf{题目}：一个10Mbps采用CSMA/CD基带总线的以太网，总线长度为1000米，信号传播速度为200米/微秒，假如网络的两个节点位于总线两端，试问：(1)两节点间的信号传播延迟是多少？(2)最多经过多长时间才能检测到冲突？(3)最小帧长度是多少？
\tcblower
(1) 信号传播延迟 = $\frac{1000\text{米}}{200\text{米/微秒}} = \textbf{5微秒}$

(2) 最大的冲突检测时间 = 2 × 信号传播延迟 = $2 \times 5 = \textbf{10微秒}$

(3) 最小帧长度 = 发送速率 × 冲突检测时间 = $10\text{Mbps} \times 10\mu s = \textbf{100比特}$
\end{examplebox}

\begin{examplebox}[计算题3.10 曼彻斯特编码]
\textbf{题目}：画出0001110101的曼彻斯特编码和差分曼彻斯特编码（假定起始状态为低）。
\tcblower
\textbf{曼彻斯特编码规则}：
\begin{itemize}[itemsep=0pt]
    \item 比特1：码元中间从高到低跳变
    \item 比特0：码元中间从低到高跳变
\end{itemize}

\textbf{差分曼彻斯特编码规则}：
\begin{itemize}[itemsep=0pt]
    \item 码元中间始终有跳变（用于同步）
    \item 比特0：码元开始处有跳变
    \item 比特1：码元开始处无跳变
\end{itemize}

数据：0 0 0 1 1 1 0 1 0 1

（具体波形图需要根据规则绘制）
\end{examplebox}

\begin{examplebox}[简答题：PPP协议工作状态]
\textbf{题目}：简述PPP协议的工作状态有哪几种？每一种连接解决什么问题？
\tcblower
\begin{enumerate}[itemsep=0pt]
    \item \textbf{链路静止状态}：不存在物理层的连接
    \item \textbf{链路建立状态}：建立链路层的LCP连接
    \item \textbf{鉴别状态}：身份验证
    \item \textbf{网络层协议状态}：进行网络层配置
\end{enumerate}
\end{examplebox}

\begin{examplebox}[计算题3.11 信道利用率]
\textbf{题目}：长度为100字节的应用层数据，从上至下依次传递给传输层、网络层和数据链路层（以太网），已知以太网帧头部为X字节，最后从网络接口传送出至少多少位比特序列？信道有效利用率为多少？
\tcblower
假设TCP首部20字节，IP首部20字节，以太网帧头部18字节（含FCS）

总长度 = 100 + 20 + 20 + 18 = 158字节 = \textbf{1264位}

信道有效利用率 = $\frac{100}{158}$ = \textbf{63.3\%}

（若X=18，则为$\frac{100}{140+X} = \frac{100}{158}$）
\end{examplebox}

\begin{examplebox}[计算题3.12 滑动窗口效率]
\textbf{题目}：数据链路层采用后退N帧协议，发送方已发送了编号为0-7的帧，当计时器超时时，只收到0、1、2号帧的确认。发送方需要重发哪些帧？
\tcblower
后退N帧协议采用累积确认，收到2号帧的确认表示0、1、2号帧都已正确接收。

超时时，需要重发从3号帧开始的所有已发送帧。

\textbf{需要重发}：3、4、5、6、7号帧（共5帧）
\end{examplebox}

\begin{examplebox}[计算题3.13 CRC校验]
\textbf{题目}：数据为1101011011，生成多项式G(x)=$x^4+x+1$，求CRC校验码和发送的比特序列。
\tcblower
\textbf{步骤}：
\begin{enumerate}[itemsep=0pt]
    \item 生成多项式对应二进制：$x^4+x+1$ = 10011
    \item 数据后补4个0：11010110110000
    \item 模2除法：
\end{enumerate}

\begin{verbatim}
    11010110110000 ÷ 10011
    余数：1110
\end{verbatim}

\textbf{CRC校验码}：1110

\textbf{发送序列}：11010110111110
\end{examplebox}

\begin{examplebox}[选择题（补充）]
\textbf{8.} 1-坚持CSMA方式中，"坚持"的含义是（\textbf{A}）\\
A. 当侦听到总线忙时，坚持侦听直到总线空闲为止\\
B. 当侦听到总线忙时，坚持发送其数据帧\\
C. 发送数据后，坚持侦听总线看是否发生冲突\\
D. 当发生冲突时，坚持发送数据

\textbf{9.} 数据链路层的协议数据单元为（\textbf{C}）\\
A. 比特 \quad B. 报文分组 \quad C. 帧 \quad D. 报文

\textbf{10.} 滑动窗口协议中，接收窗口保存的是（\textbf{B}）\\
A. 可发送的帧序号 \quad B. 可接收的帧序号\\
C. 不可发送的帧序号 \quad D. 不可接收的帧序号

\textbf{11.} CSMA/CD的访问控制方式是IEEE 802.3标准中制定的，其中CD是指（\textbf{A}）\\
A. 冲突检测 \quad B. 码分 \quad C. 激光唱盘 \quad D. 呼叫设备

\textbf{12.} 以太网采用的退避算法是（\textbf{A}）\\
A. 二进制指数截断退避算法 \quad B. 随机退避算法\\
C. 线性退避算法 \quad D. 定时退避算法
\end{examplebox}

\begin{examplebox}[综合题3.14 以太网最小帧长]
\textbf{题目}：以太网规定最小帧长为64字节，若数据传输速率为100Mbps，信号传播速度为$2\times10^8$m/s，求以太网的最大段长度。
\tcblower
最小帧长 = 64字节 = 512比特

发送最小帧的时间 = $\frac{512}{100\times10^6}$ = 5.12$\mu$s

争用期 = 发送最小帧时间 = 5.12$\mu$s

单向传播时延 = $\frac{\text{争用期}}{2}$ = 2.56$\mu$s

最大段长度 = 传播速度 × 单向传播时延 = $2\times10^8 \times 2.56\times10^{-6}$ = \textbf{512米}
\end{examplebox}
